\session{9}{10/31}{2限}{3}{可視化と探索的分析}

\goals{
  \item 伝えたい内容に応じてグラフの種類を選べる。
  \item 誤解を招く可視化の典型を見抜き、避けられる。
  \item 探索的データ分析（EDA）の姿勢を理解し、AIにグラフを作らせて人が解釈する役割分担を体験する。
}

\guidesec{解説}

\topic{グラフの選び方}
グラフは「何を伝えたいか」で選びます。見た目の好みや目新しさで選ぶと、伝わらないグラフになります。
\par\noindent\begin{gtabh}{colspec={Q[l,wd=24mm] Q[l,wd=34mm] X[l]}}
  伝えたいこと & 適したグラフ & 注意 \\
  量の比較 & 棒グラフ & 縦軸は0から始める。項目の並び順に意味を持たせる（大きい順、時系列順）。 \\
  時間の推移 & 折れ線グラフ & 間隔が等しい時間軸に使う。系列が多すぎると読めない（5本程度まで）。 \\
  構成比 & 積み上げ棒グラフ、帯グラフ & 円グラフは項目が多いと比較しにくい。構成比の変化は帯グラフで。 \\
  2変数の関係 & 散布図 & 相関の有無、傾向、外れ値を見る。相関は因果ではない（第10回）。 \\
  分布 & ヒストグラム、箱ひげ図 & 平均だけでは分からない偏りや外れ値を見る。 \\
  地理的な分布 & 地図（コロプレス図） & 面積の大きい地域が目立ちすぎる点に注意。 \\
\end{gtabh}
加えて、タイトル（何のグラフか）、軸のラベルと単位、凡例、出典を必ず示します。これらがないグラフは、読み手が正しく解釈できません。

\topic{誤解を招く可視化}
同じデータでも、描き方で印象は大きく変わります。意図的でなくても、次のような描き方は読み手を誤解させます。
\begin{itemize}
  \item \textbf{縦軸の原点が0でない棒グラフ}　わずかな差が大きな差に見える。
  \item \textbf{軸の範囲や期間の恣意的な切り出し}　都合のよい期間だけを示して、傾向を誇張する。
  \item \textbf{二重軸}　左右の軸の尺度を調整することで、無関係な2系列が連動しているように見せられる。
  \item \textbf{3D表示や装飾}　奥行きや遠近で量の比較が歪む。
  \item \textbf{面積や体積の誤用}　2倍の値を2倍の高さと2倍の幅で描くと、面積は4倍になる。
  \item \textbf{累積値の使用}　累積のグラフは常に右肩上がりになり、減速が見えにくい。
\end{itemize}
AIが作ったグラフにもこれらの問題は起こります。特に、AIは「見栄えのよい」グラフを作る傾向があり、原点や軸の範囲を自動で調整することがあります。

\topic{探索的データ分析の姿勢}
探索的データ分析（EDA：Exploratory Data Analysis）は、統計学者テューキーが提唱した考え方で、仮説を固める前に、
データを様々な角度から眺めて、分布・傾向・異常値・関係を把握する作業です。結論を先に決めて、それに合う図だけを作るのは、
確証バイアス（第11回）の典型であり、EDA の逆です。EDA では次のことを行います。
\begin{itemize}
  \item 各変数の分布を見る（ヒストグラム、要約統計量：平均、中央値、最小・最大）。
  \item 欠損と外れ値を確認し、原因を考える（入力ミスか、本当に特異な事例か）。
  \item 変数同士の関係を見る（散布図、相関）。
  \item 部分集団（店舗別、顧客層別、期間別）に分けて見る。全体の傾向と部分の傾向が異なることがある（シンプソンのパラドックス）。
  \item 気づいたことと疑問を記録し、仮説の候補にする。
\end{itemize}

\topic{AIにグラフを作らせ、人が解釈する}
生成AIは、データを渡すと数秒でグラフを作ります。複数の種類のグラフを作らせて比較することも容易です。
この速さを活かすと、EDA の「様々な角度から眺める」作業が効率的になります。一方で、どのグラフが目的に合っているか、
誤解を招いていないか、グラフから何が言えて何が言えないかを判断するのは人の役割です。
グラフは「データを見る道具」であって「結論」ではない、という点を忘れないでください。

\topic{キーワード}
\par\noindent\begin{keywords}{}
  棒・折れ線・散布図・ヒストグラム & 量の比較、時間の推移、2変数の関係、分布を示す基本のグラフ。 \\
  誤解を招く可視化 & 原点の切り詰め、期間の恣意的な選択、二重軸、3D表示など。 \\
  探索的データ分析（EDA） & 仮説を固める前に、データを様々な角度から眺めて特徴を把握する作業。 \\
  要約統計量 & 平均、中央値、最小・最大、標準偏差など、分布を要約する数値。 \\
  シンプソンのパラドックス & 全体で見た傾向と、部分集団ごとに見た傾向が逆転する現象。 \\
\end{keywords}

\guidesec{演習課題}

\exercise{1}{AI演習：同じデータから複数のグラフを作らせ、最も伝わる表現を選ぶ}
\instr{第8回で整えた公開統計データ（または次の表）を使い、次のプロンプトで3種類以上のグラフを作らせてください。各グラフについて「伝わること」「伝わらないこと」「誤解を招く点」を書き、「2026年の商品Bの不調を経営会議で伝える」という目的に最も合うグラフを選んで理由を述べてください。}
\par\noindent\begin{gtab}{colspec={Q[l,wd=18mm] Q[r,wd=20mm] Q[r,wd=20mm] Q[r,wd=20mm] Q[r,wd=20mm]}}
  年 & 商品A & 商品B & 商品C & 合計 \\
  2022 & 420 & 380 & 150 & 950 \\
  2023 & 450 & 390 & 190 & 1,030 \\
  2024 & 470 & 370 & 240 & 1,080 \\
  2025 & 480 & 340 & 300 & 1,120 \\
  2026 & 490 & 290 & 360 & 1,140 \\
\end{gtab}
\begin{promptbox}[複数のグラフの作成]
次の表は、ある会社の商品別の売上（百万円）です。この表から、(1) 折れ線グラフ、(2) 積み上げ棒グラフ、(3) 100\%積み上げ（帯）グラフ、(4) 商品Bだけの棒グラフを作ってください。各グラフには、タイトル・軸ラベル・単位・凡例を付け、縦軸は0から始めてください。それぞれのグラフが「伝えやすいこと」と「伝えにくいこと」を一言ずつ添えてください。
（表）
\end{promptbox}

\exercise{2}{誤解を招くグラフを直す}
\instr{上の表の「合計」を、縦軸を 900 から 1,200 の範囲にした棒グラフで示すと、読み手はどう受け取るか説明してください。そのうえで、「売上は堅調に伸びているが、伸びは年々鈍化している」という事実を誤解なく伝えるグラフを設計してください（グラフの種類、軸、補助線、注記）。}

\exercise{3}{EDA の記録}
\instr{上の表について、EDA の手順に沿って「気づいたこと」を5つ、「次に確かめたい疑問」を3つ書いてください。疑問は「どのデータがあれば答えられるか」も添えます。}

\guidesec{模範解答}

\begin{answer}{演習1}
\par\noindent\begin{gtabh}{colspec={Q[l,wd=28mm] X[l] X[l] X[l]}}
  グラフ & 伝わること & 伝わらないこと & 誤解を招く点 \\
  折れ線（3商品） & 商品Bの減少と商品Cの増加の対比。交差する年。 & 合計の規模。 & 縦軸の範囲を狭くすると商品Aの微増が急増に見える。 \\
  積み上げ棒 & 合計の伸びと、その中での商品の入れ替わり。 & 商品Bの減少（上に積まれた部分の変化は読みにくい）。 & 中段の系列は高さの比較が難しい。 \\
  帯（100\%） & 構成比の変化。商品Bの比率の低下が明確。 & 売上の絶対額と合計の伸び。 & 「合計が減っている」と誤解されうる。 \\
  商品Bのみの棒 & 商品Bの減少が最も直接的に伝わる。 & 他の商品との関係、会社全体の状況。 & 単独で示すと、全体が不調という印象を与えかねない。 \\
\end{gtabh}
選択の例：経営会議で「商品Bの不調」を伝えるなら、折れ線グラフ（3商品）を主に使い、合計は注記で示す。
商品Bが減っているのに合計が伸びている事実を同時に見せることで、「商品Cへの移行が起きているのか、商品Bの固有の問題か」という
次の問いにつながる。商品Bだけの棒グラフは、文脈を切り離してしまうため補助的に使う。
\end{answer}

\begin{answer}{演習2}
縦軸を 900〜1,200 にすると、2022年の棒はほとんど高さがなく、2026年の棒は4倍以上の高さになり、「売上が数倍に伸びた」という印象を与えます。
実際の伸びは5年で約20\%です。

設計例：縦軸を0から始めた棒グラフで合計を示し、各年の前年比伸び率（8.4\%、4.9\%、3.7\%、1.8\%）を棒の上に数値で添えるか、
右軸を使わずに別の小さな折れ線グラフとして並べる。注記に「伸び率は前年比」と明記する。
こうすると、「伸びている」（棒の高さ）と「鈍化している」（伸び率の低下）が、どちらも誇張なく伝わります。
\end{answer}

\begin{answer}{演習3}
気づいたこと：(1) 合計は毎年増えているが伸び率は低下。(2) 商品Bは2023年をピークに減少し、2026年は2022年比で約24\%減。
(3) 商品Cは5年で2.4倍。(4) 商品Aは安定して微増。(5) 商品BとCの合計は 530、580、610、640、650 と増えており、BからCへの置き換えが起きている可能性。

疑問：(1) 商品BとCは同じ顧客層が買っているか（顧客別の購買データがあれば答えられる）。(2) 商品Bの減少は価格の変更や競合の影響か
（価格の履歴と競合の動向）。(3) 月別に見ると減少は特定の時期に起きているか（月別売上データ）。
\end{answer}

\guidesec{出席確認クイズの解説}
講義サイトの出席確認ページ（第9回）の5問です。
% 出席確認クイズ 第09回　可視化と探索的分析
%   attendance/attendance09.html から生成（解説は本ガイド用に加筆）。

\quizq{1}{時系列の推移を示すのに一般的に適したグラフはどれでしょうか?}%
  {ヒストグラム}%
  {円グラフ}%
  {散布図}%
  {\correct{折れ線グラフ}}%
  {正解は (d)。時間に沿った変化を示すには折れ線グラフが適しています。(a) ヒストグラムは分布、(b) 円グラフは構成比、(c) 散布図は2変数の関係を見るためのグラフです。グラフの種類は「何を伝えたいか」で選びます。}

\quizq{2}{2つの量的変数の関係を見るのに適したグラフはどれでしょうか?}%
  {棒グラフ}%
  {帯グラフ}%
  {\correct{散布図}}%
  {円グラフ}%
  {正解は (c)。散布図は、2つの量的変数の関係（相関の有無や傾向、外れ値）を視覚的に確認するのに適しています。(a) 棒グラフは量の比較、(b) 帯グラフと (d) 円グラフは構成比を示すためのものです。}

\quizq{3}{誤解を招く可視化の例として適切なものはどれでしょうか?}%
  {凡例を付ける}%
  {軸にラベルを付ける}%
  {単位を明記する}%
  {\correct{縦軸の原点を0にせず、差を誇張して見せる}}%
  {正解は (d)。縦軸の原点を0以外にした棒グラフは、わずかな差が大きな差のように見え、読み手を誤解させます。(a)(b)(c) の凡例・軸ラベル・単位の明記は、誤解を防ぐための良い実践です。ほかにも3D表示や二重軸、都合のよい期間の切り出しが誤解を招く例です。}

\quizq{4}{「探索的データ分析(EDA)」の説明として適切なものはどれでしょうか?}%
  {分析を外部に丸投げする}%
  {結論を先に決めてから、それに合う図だけを作る}%
  {\correct{仮説を固める前に、データを様々な角度から眺めて特徴を把握する}}%
  {データを見ずに報告書を書く}%
  {正解は (c)。探索的データ分析（EDA）は、本格的な分析や仮説の確定の前に、データの分布・傾向・異常値などを様々な角度から確認する作業です。(b) は結論を先に決めて都合のよい図だけを作る、確証バイアス（第11回）の典型例です。}

\quizq{5}{生成AIにグラフを作らせたときの人の役割として適切なものはどれでしょうか?}%
  {見た目が派手なものを選ぶ}%
  {グラフの種類は問わない}%
  {AIが作ったので確認は不要}%
  {\correct{目的に合っているか、誤解を招かないかを判断する}}%
  {正解は (d)。AIはグラフを素早く作れますが、目的に合っているか、誤解を招かないかの判断は人の役割です。(a) 派手さは判断基準ではなく、(c) AIが作ったグラフでも軸や単位の誤りはあります。AIに作らせ、人が解釈する、という役割分担が本講義の基本です。}

